%% LyX 2.4.4 created this file.  For more info, see https://www.lyx.org/.
%% Do not edit unless you really know what you are doing.
\documentclass[english]{article}
\usepackage[T1]{fontenc}
\usepackage[latin9]{inputenc}
\usepackage{enumitem}
\usepackage{amsmath}
\usepackage{amsthm}
\usepackage{amssymb}
\usepackage{geometry}
\geometry{verbose,tmargin=1in,bmargin=1in,lmargin=1in,rmargin=1in}

\makeatletter
%%%%%%%%%%%%%%%%%%%%%%%%%%%%%% Textclass specific LaTeX commands.
\numberwithin{equation}{section}
\numberwithin{figure}{section}
\newlength{\lyxlabelwidth}      % auxiliary length 
\theoremstyle{definition}
\newtheorem*{defn*}{\protect\definitionname}

%%%%%%%%%%%%%%%%%%%%%%%%%%%%%% User specified LaTeX commands.
\usepackage{fancyhdr}
\usepackage{lastpage}
\pagestyle{fancy}
\usepackage{enumitem}
\usepackage{ifthen}
\everymath=\expandafter{\the\everymath\displaystyle}
%\DeclareMathSizes{10}{10}{10}{10}
\usepackage{multicol}

\usepackage{tikz}
\usetikzlibrary{patterns}
\usetikzlibrary{plotmarks}
\usepackage{pgfplots}
\pgfplotsset{compat=newest}

\usepackage{calc}
\usepackage[nomessages]{fp}

\usepackage{datetime}
% Added by lyx2lyx
\setlength{\parskip}{\smallskipamount}
\setlength{\parindent}{0pt}

\makeatother

\usepackage{babel}
\providecommand{\definitionname}{Definition}

\begin{document}
\renewcommand{\author}{Dr.\,James Ashe}

\newcommand{\classNum}{335}

\newcommand{\classSec}{A}

\newcommand{\nameType}{Name: \hspace{-4cm}}

\newcommand{\examType}{Homework 5}

\newcommand{\Date}{Due date: \formatdate{26}{4}{2013}} %%\AdvanceDate[12] \today

%\newdate{my_date}{\day}{\month}{\year}%   
%\AdvanceDate[-5] 
%\displaydate{my_date}

%%First page's header%%%%%%%%%%%%%%%%%%%%%%
\fancypagestyle{firststyle} %%header settings for first pagr
{
	\fancyhead[L]{MTH \classNum{}(\classSec{}) \author{}\\
	\textbf{\examType{}} \Date{}}
	\fancyhead[C]{\textbf{\nameType}}
	\setlength{\headsep}{14pt}
}
%%%%%%%%%%%%%%%%%%%%%%%%%%%%%%%%%%%%%%%%%%

%%Next page's header%%%%%%%%%%%%%%%%%%%%%%%
%\setlist{nolistsep} %eliminates space between list items
\lhead{MTH \classNum{}(\classSec{})} 
\chead{\textbf{\nameType}} 
\cfoot{\thepage{}~of~\pageref{LastPage}} 
\setlength{\headsep}{5pt} 
%%%%%%%%%%%%%%%%%%%%%%%%%%%%%%%%%%%%%%%%%%%

%%Various settings; see the preamble for other settings%%%%%
%%\setenumerate[0]{leftmargin=12pt,itemindent=0pt,labelwidth=0pt}
\renewcommand{\labelenumi}{\textbf{\arabic{enumi}.}}
\renewcommand{\labelenumii}{\textbf{\alph{enumii}.}}
\thispagestyle{firststyle}
%%%%%%%%%%%%%%%%%%%%%%%%%%%%%%%%%%%%%%%%%%

\textbf{Homework Problems. }\emph{Solutions should be written/typed
in numerical order and clearly labeled. Unsolved problems should be
included but left blank.}
\begin{enumerate}
\item Let $\left(\mathbb{Q},+\right)$ be the group of rational numbers
under addition. 

\begin{enumerate}
\item Is $\mathbb{Z}$ a subgroup of $\left(\mathbb{Q},+\right)$?

\begin{proof}
Yes. By the subgroup test: $\mathbb{Z}\neq\emptyset$ (since $1\in\mathbb{Z}$)
and $\mathbb{Z}\subseteq\mathbb{Q}$ (since $a\in\mathbb{Z}\Rightarrow\frac{a}{1}\in\mathbb{Q}$).
If $a,b\in\mathbb{Z}$ then $b^{-1}=-b$ so $ab^{-1}=a-b\in\mathbb{Z}$.
\\
Alternatively, simply note that $\mathbb{Z}\subseteq\mathbb{Q}$ and
that $\left(\mathbb{Z},+\right)$ is a group by a previous exercise.
\end{proof}
\item Is $\mathbb{Z}^{+}=\left\{ 0,1,2,3,\dots\right\} $ a subgroup of
$\left(\mathbb{Q},+\right)$?

\begin{proof}
No. If $a\in\mathbb{Z}^{+}$ then $aa^{-1}=0$ implies that $a^{-1}=-a$,
and $-a\notin\mathbb{Z}$ if $a\geq0$ (groups/subgroups must contain
the inverse of each element). 
\end{proof}
\end{enumerate}
\item Determine if the following are subgroups.

\begin{enumerate}
\item Is $\mathbb{Q}$ a subgroup of the group, $\left(\mathbb{C},+\right)$,
the complex numbers under addition?

\begin{proof}
Yes. By the subgroup test: If $\frac{a}{b}\in\mathbb{Q}$ for $a,b\in\mathbb{Z}$
then $\frac{a}{b}=\frac{a}{b}+0\cdot i\in\mathbb{C}$, and $\frac{1}{1}\in\mathbb{Q}\Rightarrow\mathbb{Q}\neq\emptyset$.
If $\frac{a}{b},\frac{c}{d}\in\mathbb{Q}$ for some $a,b,c,d\in\mathbb{Q}$
then $\left(\frac{c}{d}\right)^{-1}=-\frac{c}{d}$, and $\frac{a}{b}-\frac{c}{d}=\frac{ad-cb}{bd}\in\mathbb{Q}$
since $ad-cb$ and $bd\in\mathbb{Z}$.
\end{proof}
\item Is $\mathbb{Z}-\left\{ 0\right\} $ a subgroup of the group, $\left(\mathbb{Q}-\left\{ 0\right\} ,\cdot\right)$,
the non-zero rational numbers under multiplication?

\begin{proof}
No. $e=1$ so if $a\in\mathbb{Z}$ and $a\neq1$ then $a\cdot a^{-1}=1\Rightarrow a^{-1}=\frac{1}{a}\notin\mathbb{Z}$. 
\end{proof}
\end{enumerate}
\item Prove that the intersection of two subgroups $H$ and $K$ of a group
$G$ is a subgroup.

\begin{proof}
Let $H,K\leq G$ and $a,b\in H\cap K$. By the subgroup test: Since
$H$ and $K$ are subgroups $e\in H$ and $e\in K$$\Rightarrow e\in H\cap K$
where $e$ is the identity of $G$. Thus $H\cap K\neq\emptyset$.
Also $H\cap K\subseteq G$. $b\in H$ implies that $b^{-1}\in H$
since $H$ is a subgroup. So then $ab^{-1}\in H$ since $a\in H$
and both $H$ is closed. Similarly $ab^{-1}\in K$ so then $ab^{-1}\in H\cap K$,
and $H\cap K\leq G$.
\end{proof}
\item Let $H_{1},H_{2},\dots$ be subgroups of $G$. Prove that $H_{1}\cap H_{2}\cap\cdots$
is a subgroup of $G$.

\begin{proof}
$H_{1}\cap H_{2}\cap\cdots={\displaystyle \cap_{i=1}^{\infty}H_{i}}$.
By the subgroup test: $e\in\cap_{i=1}^{\infty}H_{i}$ and each $H_{i}\leq G$
so then $\cap_{i=1}^{\infty}H_{i}\neq\emptyset$ and $\cap_{i=1}^{\infty}H_{i}\subseteq G$.
If $a,b\in H_{i}$ then $b^{-1}\in H_{i}$ for each $i$ since every
$H_{i}$ is a subgroup. So then $ab^{-1}\in H_{i}$ since each $H_{i}$
is closed which implies than $ab^{-1}\in\cap_{i=1}^{\infty}H_{i}$.
\end{proof}
\item Recall that $\mathbb{Z}_{6}=\left\{ 0,1,2,3,4,5\right\} $ is a group
under modular addition. 

\begin{enumerate}
\item Show that $H=\left\{ 0,2,4\right\} $ and $K=\left\{ 0,3\right\} $
are subgroups of $\mathbb{Z}_{6}$.

\begin{proof}
$H=\left\langle 2\right\rangle $ and $K=\left\langle 3\right\rangle $,
i.e., they are both cyclic. So they are both subgroups by theorem
stated in class. (\textbf{Thm: }If $G$ is a subgroup and $a\in G$
then $\left\langle a\right\rangle \leq G$ under the operation of
$G$). 
\end{proof}
\item Show that $H\cup K$ is \emph{not }a subgroup of $\mathbb{Z}_{6}$.

\begin{proof}
$H\cup K=\left\{ 0,2,3,4\right\} $. $2,3\in H\cup K$, but $2+3=5\notin H\cup K$. 
\end{proof}
\item Is the union of two subgroups of $G$ necessarily a subgroup $G$?

\begin{proof}
No. $H$ and $K$ are both subgroups by part (a), but $H\cup K$ is
not a subgroup by part (b).
\end{proof}
\end{enumerate}
\item Let $G$ be an abelian group and $H$ be a subgroup of $G$. Prove
that $S(H)=\left\{ x|x\in G\mbox{ and }x\cdot x\in H\right\} $ is
also a subgroup of $G$.

\begin{proof}
By definition, $S(H)\subseteq G$. Also the identity is in every subgroup
so $e\in H$. Thus $e\cdot e\in S(H)$, and $S(H)\neq\emptyset$.
If $x,y\in S(H)$ then $x\cdot x\cdot\in H$ and $y\cdot y\in H\Rightarrow y^{-1}y^{-1}\in H$
since $H$ is a subgroup. So then $\left(x\cdot x\right)\left(y^{-1}y^{-1}\right)=xy^{-1}xy^{-1}$
since $G$ is abelian. Thus $xy^{-1}\in S(H)$, and $S(H)\leq G$.
\end{proof}
\item Let $G$ be a group. Prove that $G$ has a subgroup with one element.
This group is called the \emph{trivial group}.

\begin{proof}
Because $G$ is a group it has an identity, $e$. $\left\{ e\right\} \subseteq G$,
$\left\{ e\right\} \neq\emptyset$, and $e\cdot e^{-1}=e\cdot e=e\in\left\{ e\right\} $.
Thus $\left\{ e\right\} \leq G$. (Alternatively, note that $\left\{ e\right\} $
is cyclic and use the theorem as in problem 5.a).
\end{proof}
\end{enumerate}
\begin{defn*}
Let $G$ be a group. If there is a single $g\in G$ such that every
element in $G$ is equal to $g^{n}$ for some $n\in\mbox{\ensuremath{\mathbb{N}}}$
then $G$ is a \emph{cyclic }group and \emph{x }is called the \emph{generator}
of $G$.
\end{defn*}
\begin{enumerate}[resume]
\item Prove that:

\begin{enumerate}
\item $\mathbb{Z}_{n}$ is a cyclic group under modular addition for any
$n\in\mathbb{Z}$ such that $n>0$. (you can use the fact that it
is a group).

\begin{proof}
$n>0$, so then $1\in\mathbb{Z}_{n}$, and $\left\langle 1\right\rangle =\left\{ 1,2,3,\dots,n-1,0\right\} =\mathbb{Z}_{n}$.
Thus $\mathbb{Z}_{n}$ is cyclic generated by $1$.
\end{proof}
\item $\mathbb{Z}$ is a cyclic group. 

\begin{proof}
We have already shown that $\mathbb{Z}$ is a group. $1\in\mathbb{Z}$,
and $\left\langle 1\right\rangle =\left\{ \dots,-2,-1,0,1,2,\dots\right\} $.
($1^{-2}=-2$; As stated in class, all cyclic groups of the same order
are \emph{isomorphic}. So $\mathbb{Z}$ is the only infinite cyclic
group up to isomorphism).
\end{proof}
\end{enumerate}
\item Find all generators of 

\begin{enumerate}
\item $\mathbb{Z}_{6}$, $\mathbb{Z}_{8}$, and $\mathbb{Z}_{10}$ (note:
$\mathbb{Z}_{n}$ is often notated $Z_{n}$).

\begin{enumerate}
\item [$\mathbb{Z}_{6}$]$\left\langle 1\right\rangle =\left\{ 1,2,3,4,5,0\right\} $(note:
$1$ is a generator of $\mathbb{Z}_{n}$ for any $n$) and $\left\langle 5\right\rangle =\left\{ 5,4,3,2,1,0\right\} $
so $\left\langle 1\right\rangle =\left\langle 5\right\rangle =\mathbb{Z}_{6}$.
\item [$\mathbb{Z}_{8}$]$\left\langle 1\right\rangle $, $\left\langle 3\right\rangle =\left\{ 3,6,1,4,7,2,5,0\right\} $,
$\left\langle 5\right\rangle =\left\{ 5,2,7,4,1,6,3,0\right\} $,
$\left\langle 7\right\rangle =\left\{ 7,6,5,4,3,2,1,0\right\} $.
\item [$\mathbb{Z}_{10}$]Similarly to above we find that the generators
are:$\left\langle 1\right\rangle $, $\left\langle 3\right\rangle $,$\left\langle 7\right\rangle $,
and $\left\langle 9\right\rangle $
\end{enumerate}
\end{enumerate}
\end{enumerate}
\begin{defn*}
Let $X=\left\{ 1,2,\dots,n\right\} $. The \emph{symmetric group on
X, $S_{n}$, }is the collection of all bijections $\sigma:X\rightarrow X$.
That is, it is the set of all functions that rearrange the elements
of $X$. ($S_{n}$ is also called the \emph{permutation group of the
set X}).
\end{defn*}
\begin{enumerate}[resume]
\item Let $S_{3}$ be the symmetric group of permutations on three elements.
Show that $S_{3}$ is:

\begin{enumerate}
\item Non-abelian. 

\begin{proof}
Let $\sigma_{1}=\left(\begin{array}{ccc}
1 & 2 & 3\\
2 & 3 & 1
\end{array}\right)$ and $\sigma_{2}=\left(\begin{array}{ccc}
1 & 2 & 3\\
1 & 3 & 2
\end{array}\right)$. $\sigma_{1}\sigma_{2}=\left(\begin{array}{ccc}
1 & 2 & 3\\
3 & 2 & 1
\end{array}\right)\neq\left(\begin{array}{ccc}
1 & 2 & 3\\
2 & 1 & 3
\end{array}\right)=\sigma_{2}\sigma_{1}$.  
\end{proof}
\item Has a cyclic subgroup of order $3$. 

\begin{proof}
$\sigma_{1}^{1}=\left(\begin{array}{ccc}
1 & 2 & 3\\
2 & 3 & 1
\end{array}\right)$, $\sigma_{1}^{2}=\left(\begin{array}{ccc}
1 & 2 & 3\\
2 & 3 & 1
\end{array}\right)\left(\begin{array}{ccc}
1 & 2 & 3\\
2 & 3 & 1
\end{array}\right)=\left(\begin{array}{ccc}
1 & 2 & 3\\
3 & 1 & 2
\end{array}\right)$, and $\sigma_{1}^{3}=\sigma_{1}^{2}\sigma_{1}=\left(\begin{array}{ccc}
1 & 2 & 3\\
3 & 1 & 2
\end{array}\right)\left(\begin{array}{ccc}
1 & 2 & 3\\
2 & 3 & 1
\end{array}\right)=\left(\begin{array}{ccc}
1 & 2 & 3\\
1 & 2 & 3
\end{array}\right)=e$. So then $\left|\sigma_{1}\right|=3$, and we know that $\left\langle \sigma_{1}\right\rangle \leq S_{3}$.\vfill
\end{proof}
\end{enumerate}
\end{enumerate}

\end{document}
